\section{Compline, Examination of Conscience, and Received Prayer}

The Church does not wait for a grave medical prognosis to teach her children how to die. Every evening, as daylight fades and the body prepares to surrender consciousness to sleep---that nightly rehearsal of the grave---the Roman liturgy places Compline upon the lips of the faithful. Compline is the Church's daily school of holy dying.

In the classical pre-1911 Roman Breviary arrangement preserved in the 1889 \emph{Manual of Prayers}, the hour opens with St.\ Peter's watchword of vigilance: ``Be sober and watch: because your adversary the devil, as a roaring lion, goeth about seeking whom he may devour'' (1 Pet.\ 5:8). After examining the day and confessing sin in the \emph{Confiteor}, the believer prays four unchanging Psalms every night before sleep.\endnote{\emph{Manual of Prayers for the Use of the Catholic Laity}, prepared and published by order of the Third Plenary Council of Baltimore (New York: Catholic Publication Society, 1889), Evening Prayers (Compline), pp.\ 69--75 (including Cardinal Newman's translation of \emph{Te lucis ante terminum}, p.\ 73), reflecting the classical pre-1911 Roman Breviary Psalter.} Psalm 4 bids the heart inspect itself in stillness---``the things ye say in your hearts, be sorry for upon your beds''---and lie down without dread: ``In peace in the self-same: I will sleep, and I will rest. For Thou, O Lord, alone: hast established me in hope'' (Ps.\ 4:5, 9--10). Psalm 30 leads directly to Calvary: ``Into Thy hands I commend my spirit: Thou hast redeemed me, O Lord, God of truth'' (Ps.\ 30:6). Psalm 90 wraps the night in God's protection and the guard of His angels (Ps.\ 90:1, 11), and Psalm 133 lifts the hands of those who stand by night in the courts of the Lord's house. Subsequent Roman reforms varied this psalmody across the week while keeping its night-watch core: Pope St.\ Pius X's 1911 Psalter reform (\emph{Divino afflatu}), carried forward in the 1962 Roman Breviary, redistributed the Compline psalms across the weekdays---appointing a three-psalm Sunday Compline (Psalms 4, 90, and 133) and moving Psalm 30 to Wednesday---while the 1971 \emph{Liturgia Horarum}, promulgated by Pope St.\ Paul VI's \emph{Laudis canticum} (1970), likewise varies the nightly psalms across the week while retaining \emph{In manus tuas, Domine} and the \emph{Nunc dimittis} every evening.\endnote{Pope St.\ Pius X, Apostolic Constitution \emph{Divino afflatu} (1 November 1911), rearranging the Roman Breviary Psalter (maintained in the 1962 Roman Breviary); Pope St.\ Paul VI, Apostolic Constitution \emph{Laudis canticum} (1 November 1970), promulgating the 1971 \emph{Liturgia Horarum}.}

Then rises the ancient evening hymn \emph{Te lucis ante terminum}---``Now that the daylight dies away, / By all Thy grace and love, / Thee, Maker of the world, we pray / To watch our bed above''---asking that night's phantoms flee and our bodies be kept pure as shrines. After the Little Chapter, the Short Responsory repeats Christ's dying cry on the Cross, \emph{In manus tuas, Domine, commendo spiritum meum}. Then follows the Canticle of aged Simeon, the \emph{Nunc dimittis}, framed by an antiphon uniting nightly rest with final sleep in Christ: \emph{Salva nos, Domine, vigilantes, custodi nos dormientes: ut vigilemus cum Christo, et requiescamus in pace}---``Save us, O Lord, watching, guard us sleeping: that we may watch with Christ, and may rest in peace.'' In the classical Compline sequence printed in the 1889 \emph{Manual of Prayers}, after the intervening \emph{Preces}---including the Lord's Prayer and the Creed ending aloud on ``The resurrection of the body. And the life everlasting. Amen''---the Collect \emph{Visita, quaesumus, Domine} begs God to visit the house, drive far every snare of the enemy, and set His holy angels to dwell within it.

Whoever prays Compline with attention practices dying three hundred and sixty-five times a year, closing his eyes each night not in anxiety or distraction, but by surrendering his spirit deliberately into the hands of the Redeemer.

\subsection*{Examination of Conscience}

Scripture counsels every mortal believer: ``In all thy works remember thy last end, and thou shalt never sin'' (Ecclus.\ 7:40). Contemplating death bears fruit only when it turns inward into honest self-examination before God:

\begin{itemize}
\item \textit{Readiness for an unannounced hour.} If God required my soul tonight (Luke 12:20), in what state would He find me? Am I carrying any unconfessed mortal sin on my conscience, postponing the Sacrament of Penance as though tomorrow belonged to me?
\item \textit{Attachment to passing things and unfinished justice.} What earthly possession, comfort, status, or approval have I clung to as though it were permanent? Have I left debts unpaid, restitution unmade, or lawful provision for my family neglected out of avoidance?
\item \textit{Unforgiveness and broken charity.} Do I harbor resentment or refuse forgiveness to anyone who has injured me? Have I failed to seek reconciliation with someone I have wronged, letting the sun go down repeatedly upon my anger?
\item \textit{Patience under frailty and fidelity to the last sacraments.} When sickness, aging, or loss enters my life, do I murmur against Providence or unite my weakness to the Passion of Christ? Have I made sure that, if grave illness strikes me or someone in my care, a priest will be summoned promptly for Confession, Anointing of the Sick, and Holy Viaticum before consciousness fades?
\item \textit{Piety toward the dead and nightly surrender.} Do I pray and have Masses offered for my deceased parents, kindred, friends, benefactors, and the forgotten dead? Do I close each day by examining my conscience and commending my spirit into the hands of God?
\end{itemize}

\subsection*{Received Prayers}

At nightfall, in sickness, and throughout the days of health, the Church places tested words of commendation and trust upon our lips.

\subsubsection*{From Nightly Compline}

\begin{quote}
\textit{Short Responsory.}

Into Thy hands, O Lord, I commend my spirit.\\
\textit{R.} Into Thy hands, O Lord, I commend my spirit.\\
\textit{V.} Thou hast redeemed us, O Lord, God of truth.\\
\textit{R.} I commend my spirit.\\
\textit{V.} Glory be to the Father, and to the Son, and to the Holy Ghost.\\
\textit{R.} Into Thy hands, O Lord, I commend my spirit.

\textit{V.} Keep us, O Lord, as the apple of Thine eye.\\
\textit{R.} Protect us under the shadow of Thy wings.

\textit{The Canticle of Simeon.}

Now, O Lord, lettest Thou Thy servant depart in peace: according to Thy word.\\
For mine eyes have seen: Thy salvation.\\
Which Thou hast prepared: in sight of all nations.\\
Light to enlighten the gentiles: and glory of Thy people Israel.\\
Glory be to the Father, etc.

\textit{Ant.} Save us, O Lord, watching, guard us sleeping: that we may watch with Christ, and may rest in peace.

[\ldots]

\textit{Collect.}

Visit, we beseech Thee, O Lord, this habitation, and drive far from it all snares of the enemy: let Thy holy Angels dwell therein to keep us in peace: and may Thy blessing be upon us always. Through our Lord, etc.\\
\textit{R.} Amen.\endnote{\emph{Manual of Prayers for the Use of the Catholic Laity} (1889), Evening Prayers (Compline), pp.\ 73--75 (omitting the \emph{Preces}---\emph{Kyrie}, Lord's Prayer, Apostles' Creed, and versicles---that stand between the \emph{Nunc dimittis} antiphon and the Collect).}
\end{quote}

\subsubsection*{The \textit{Anima Christi}}

\begin{quote}
Soul of Christ, sanctify me!\\
Body of Christ, save me!\\
Blood of Christ, inebriate me!\\
Water from the side of Christ, wash me!\\
Passion of Christ, strengthen me!\\
O Good Jesus, hear me!\\
Within Thy Wounds hide me!\\
Suffer me not to be separated from Thee!\\
From the malicious enemy defend me!\\
In the hour of my death, call me,\\
And bid me come to Thee;\\
That with Thy Saints I may praise Thee\\
For ever and ever. Amen.\endnote{\emph{Manual of Prayers for the Use of the Catholic Laity} (1889), Devotions for Mass, p.\ 103.}
\end{quote}

\subsubsection*{A Prayer in Preparation for Death (St.\ Alphonsus Liguori)}

\begin{quote}
Then, at death all shall be at an end for me. I shall then find only the little I have done for you, O my God. And what do I wait for? Do I wait till death come and find me as miserable and as defiled with sin as I am at present? Were I now called to eternity, I should die with great disquietude on account of my past sins. No, my Jesus; I will not die so discontented. I thank you for having given me time to weep over my iniquities and to love you. I wish to begin from this moment. I am sorry from the bottom of my heart for having offended you, O Sovereign Good, and I love you above all things,---I love you more than my life. My Jesus, I give myself entirely to you. From this moment I embrace you and unite you to my heart. I now consign my soul to you. ``Into thy hands I commend my spirit.'' I will not wait to give it to you when that \textit{proficiscere} shall announce my departure from this world. I will not wait till then to ask you to save me. \textit{Jesu sis mihi Jesus.} My Saviour, save me now by granting me pardon and the grace of your holy love. Who knows but this consideration which I have read, shall be the last call which you will give me and the last mercy which you will show me? Extend your hand, O my love, and deliver me from the mire of my tepidity. Give me fervour, and make me do with great love all that you demand of me. Eternal Father, for the love of Jesus Christ, give me holy perseverance and the grace to love you, and to love you ardently during the remainder of my life. O Mary, through the love which you bear to your Jesus, obtain for me two graces,---perseverance and love.\endnote{St.\ Alphonsus Maria de Liguori, \emph{Preparation for Death; or, Considerations on the Eternal Maxims}, trans.\ a Catholic clergyman (Dublin: James Duffy, 1843), Fourth Consideration, Third Point, Affections and Prayers, pp.\ 68--69.}
\end{quote}
